\documentclass[11pt]{article}
\usepackage[margin=1.15in]{geometry}
\usepackage{hyperref}
\hypersetup{colorlinks=true, urlcolor=blue, linkcolor=black}
\setlength{\parindent}{0pt}
\setlength{\parskip}{7pt}
\pagestyle{empty}
\begin{document}

\hfill 17 August 2026

\bigskip
Dr.\ Nur Zincir-Heywood\\
Associate Editor-in-Chief\\
IEEE Transactions on Network and Service Management

\bigskip
Dear Dr.\ Zincir-Heywood,

Thank you for the decision on TNSM-2026-11591, ``Learning to Route Under Load: An Inductive Graph
Network that Amortizes Congestion-Aware Traffic Engineering for LEO Mega-Constellations,'' and for
the invitation to submit a new manuscript. We are grateful to Associate Editor Dr.\ Tianqi Yu and
to the three reviewers, whose comments were specific enough to be acted on rather than merely
acknowledged.

We should say at the outset that this is not the same paper with corrections. Two of the reviews
asked questions we had not asked ourselves, and answering them changed our conclusions.

Reviewer 3 pointed out that our headline gap is produced by a road-traffic congestion function
evaluated where a packet link cannot operate: the excess offered to a link is lost, not delayed.
Enforcing capacity feasibility reduces that gap from $79.1\%$ measured as travel time to $30.9\%$
measured as delivered rate. Reviewer 2 pointed out that our comparators were weak, and asked for a
strong baseline at a matched computational budget. Supplying one, a congestion-blind multipath
split costing a single pass, showed that our method's advantage is decisive on the constellation it
was trained on and level with that baseline on constellations it was not.

Those two results do not support the four contributions the original manuscript claimed. Rather
than restate weakened versions of them, we have rewritten the paper as what the measurements
actually support: a study of what learned congestion prices buy in LEO routing, where a blind
baseline is already sufficient, and why the boundary falls where it does. The system model, the
simulator and the learned method are unchanged. The claims built on them are not, and the
resubmitted manuscript reports smaller numbers than the original in three places.

We also owe the Editor a correction. Following the reviewers' and your own remarks on related work,
we found arXiv:2601.21921, which trains a graph network to infer per-link congestion prices from
the constellation state in a single forward pass for LEO mega-constellations. That is the same core
idea as our original contribution, it was public five months before our submission, and we had not
cited it. It is now cited and discussed, and it is part of why the method framing is no longer the
right one for this work.

\textbf{Data and code.} The simulator, the learned model, every experiment including those added in
response to the reviewers, and the scripts that generate each reported number are available at
\url{https://github.com/haodpsut/qwgnn-leo-routing}. Every headline number in the paper is tied to
the file, column and aggregation that produce it by a claims manifest, and a script recomputes all
of them.

We hope the reviewers find the changes responsive. Where we have declined a suggestion, in
particular four suggested citations from Reviewer 1 that concern RIS and anti-jamming rather than
routing, we say so and explain why rather than adding them silently.

\bigskip
Sincerely,

\bigskip
Phuc Hao Do, PhD

\end{document}
